\section*{Acknowledgements}\label{sec:acknowledgements}

Itay Tirosh conceptualised the idea of inferring copy number alterations from single-cell expression and first applied the method that preceded inferCNV, implemented in MATLAB for the single-cell study of glioblastoma \citep{patel2014} and used in subsequent tumour studies of the Regev, Suvà and Bernstein laboratories. Brian Haas conceived of building inferCNV as a software package implementing the algorithm for general use by the community, and the package grew from there. Aviv Regev was the principal investigator of the laboratory responsible for the work. Timothy Tickle worked on the initial implementation of the inferCNV software. Christophe Georgescu further built and maintained the Bioconductor package for most of inferCNV's existence and is the largest contributor to the code. Brian Haas developed the simulation-based in-silico spike-in (h-spike) modelling used to calibrate the six-state HMM, and the random-trees subclustering, during his Ph.D. under the advisors Aviv Regev and Simon Kasif \citep{haas2021thesis}. inferCNV builds directly on HoneyBADGER by Jean Fan, Peter Kharchenko and colleagues \citep{fan2018}: the three-state expression HMM and its Kolmogorov--Smirnov calibration, the use of a Bayesian model to re-examine each HMM region, the allele HMM, and the original allele and joint expression--allele models all come from HoneyBADGER. We thank Jean Fan and Peter Kharchenko for inspiration and encouragement towards adopting features of interest from HoneyBADGER into inferCNV. Maxwell Brown integrated the Bayesian filtering steps. Yu Zhong added the allele-aware modelling, porting HoneyBADGER's allele models into inferCNV. We thank Joshua Campbell for contributions and advice, Eric Weitz for the command-line interface and Ideogram.js export, and Joshua Gould for the WDL workflow. We also thank the others who have contributed bug fixes through GitHub over the years, and the many users who reported issues.

This manuscript was prepared with the assistance of Claude Opus 5.5 (Anthropic), starting from text written by Brian Haas, Christophe Georgescu and Timothy Tickle, and structured as open-source, reproducible text with supporting analysis notebooks. It is the paper that should have been written long ago but, owing to limited time and competing priorities, never surfaced; Claude made it possible. We wrote it for inferCNV's users, and to provide an official record of its implementation and history given its wide use and impact over the years.
